\section{The claims, and what each is worth}

Eleven claims are live. Three of them are the paper: the object (C1), the unit (C3) and the price
(C4). Two are honest negatives that a reviewer will find if we do not print them (C5, C6), and one
is a scope limit that reads badly and has to be printed anyway (C8). The rest support or bound
them. Each is written below in the form it is licensed in, which is narrower than the form it is
tempting to write.

Every claim carries a line naming what would kill it. Where the register records that a novelty
search was never run, the claim is marked as blocked rather than dressed as verified.

\subsection*{C1 \quad The object: nobody prices a structural assumption}

Sensitivity analysis in causal inference is built for unmeasured confounding and for continuous
relaxations of identifying assumptions. The assumption that background knowledge carries is
discrete and structural, and it is not priced the same way. An analyst who asserts orientations,
propagates them and reads an adjustment set off the result has no number telling them how much of
the reported effect rests on those assertions. The assertions are checked for consistency with the
data and never for truth, so a claim that is consistent and false passes every check that current
practice has.

\begin{center}\footnotesize
\begin{tabular}{@{}p{2.6cm}p{13.75cm}@{}}
\toprule
rests on & the framing, not a measurement; the measurements below are what makes it a paper \\
significance & this is the one place where unverifiable human input enters an otherwise
data-driven pipeline, and the estimate at the end of the pipeline inherits it silently \\
nearest prior art & Taeb, Guo and Henckel 2025 (arXiv:2511.10625), model-oriented graph distances
over posets with an adjustment-set robustness application. This is the foundation, not the
competitor, and the paper says so. Guo and Perkovi\'c (arXiv:2010.08611) enumerate adjustment sets
in an MPDAG and never perturb the knowledge. Masten and Poirier, \emph{Quantitative Economics}
11(1):41--111, 2020, define breakdown frontiers with inference, in econometrics and not on graphs:
prior art for reporting a breakdown value with a band, outside machine-learning venues. Rosenbaum,
Cinelli and Hazlett, Ding and VanderWeele all price a continuous confounding parameter \\
falsifier & a tolerance radius centred on the analyst's own graph inside Taeb et al.\ \S6, which
nobody on this side has read in full \\
carried by & T3, the report on named analyses, written in the Rosenbaum form \\
\bottomrule
\end{tabular}
\end{center}

\begin{blocking}
\textbf{Blocked, and it blocks the abstract.} The register entry for this claim carries
\texttt{verified\_NOT\_RUN}: no novelty search has ever been run in this wording. The three sweeps
on record were pointed at latent-causal representation learning and at Taeb specifically, not at
the sensitivity-analysis lineage the claim excludes, and the framing rests on the statement that
Guo and Perkovi\'c never perturb the knowledge, from a paper nobody has read in full. Under the
submission rule a novelty claim verified more than thirty days ago cannot be reused, and this one
was never verified at all. The search has to run before the abstract goes in on 18 September, and
C1 has to be narrowed to whatever it returns.
\end{blocking}

\subsection*{C2 \quad The quantity is computable from the analyst's own inputs}

The breakdown radius is the distance from the analyst's graph to the nearest knowledge state at
which the committed adjustment set stops being valid. It is computed from the estimated CPDAG, the
claims, the treatment, the outcome and the committed set. The data-generating graph is on no path
of the computation; it is used to draw samples and, afterwards, to check whether the certificate
held.

\begin{center}\footnotesize
\begin{tabular}{@{}p{2.6cm}p{13.75cm}@{}}
\toprule
rests on & failure is upward-closed, the space is a join-semilattice with the analyst's graph as
its maximum, the nearest failure is reached by retracting only, and cover generation is set
arithmetic. Our own contribution to this is the audit: 5{,}000 sampled pairs over 469 rows across
the elicited, truthful and oracle conditions, zero violations, one-sided 95\,\% upper bound on the
violation rate of 0.06\,\% \\
significance & the reachability result rests on an anti-exchange property that is verified and not
proved, with a one-sided error that overstates robustness. The audit is the direct evidence for it,
and it is a number where the draft currently has a caveat \\
the positioning & this morning's review of the draft found the sharpest version of the point.
Taeb et al.\ prove in their Proposition 9 that the poset of MPDAGs is not graded for four or more
vertices and not semimodular, with no characterisation of the covering relation; the draft's
appendix defines a rank function and concludes that its order is graded. Both are correct, because
their counterexample varies the skeleton and this space holds it fixed. Saying so converts the
largest liability in the theory section into the clearest statement of what fixing the skeleton
buys \\
nearest prior art & the order and the distance are Taeb et al.; the specialisation to knowledge
states and the validity predicate are ours and Chris's \\
falsifier & one counterexample to upward closure halts the argument. None found in 5{,}000 draws \\
carried by & T5 makes it legible, T2 makes it safe \\
\bottomrule
\end{tabular}
\end{center}

\subsection*{C3 \quad The unit: a certificate in claims is safe, a certificate in graph revisions is not}

This is the strongest thing the measurement produced and it is the one place where our results
change what the draft reports. A certificate denominated in the analyst's own claims never
certified a set that was invalid at the truth, on any row of any condition. The same certificate
denominated in elementary graph revisions did.

\begin{center}\footnotesize
\begin{tabular}{@{}p{2.6cm}p{13.75cm}@{}}
\toprule
rests on & 3{,}014 certifiable rows under eleven knowledge conditions over 27 networks: the claim
unit over-certifies on 0, the revision unit on 39, a free rule read off the graph on 15, and the
two naive counts on 729 and 791 \\
significance & it decides the unit the paper reports in. The draft reports the radius in elementary
revisions and treats the claim count as a practitioner-facing refinement; the ledger says the
refinement is the safe one and the revision unit is not. It is also the reason the certificate has
to be redenominated rather than rescaled \\
the argument & structural, and the paper must not present it as statistical. If the committed set
is invalid at the truth and the condition asserted $w$ false claims, retracting those $w$ claims
leaves a graph the truth extends, on which the criterion fails, so the claim radius is at most $w$
whenever the search reaches depth $w$. The zero is what a sound criterion and an audit asking the
same question must return, and it was returned once the audit did \\
nearest prior art & none. No accepted paper in the record we searched computes a breakdown radius
or a null radius, so there is no precedent for the comparison and no borrowed template; the
comparison set is the instruments an analyst already has \\
falsifier & one row on which the claim certificate certifies a set invalid at the truth, under a
correct generalised adjustment criterion. Two such rows existed before the audit oracle was fixed,
and both were the oracle \\
carried by & T2, and it is the reason T2 is the recommendation \\
\bottomrule
\end{tabular}
\end{center}

\subsection*{C4 \quad The price: what the interval costs, and what it buys back}

When the claims are right, the interval that tolerates one revision costs about twice the width of
the report as committed. When one claim is wrong, the report as committed covers the true effect on
six problems in ten, and the interval covers on ninety-eight of a hundred at 85\,\% of the width of
discarding the knowledge altogether. Two wrong claims need two retractions: the interval at one
revision fails and the interval at two pays.

\begin{center}\footnotesize
\begin{tabular}{@{}p{2.6cm}p{13.75cm}@{}}
\toprule
rests on & the calibrated width table, 460 informative rows over 23 networks, with a network
cluster bootstrap; predictions P1, P2, P4 and P8 supported in four cells of four, P3 partial, P6
falsified, P7 supported \\
significance & this is the sentence a practitioner acts on, and the only one in the paper that
gives a price in the currency they report in \\
nearest prior art & the class-wide interval of Guo and Perkovi\'c is the procedure that discards
the knowledge, and it is the denominator of every cell. The omitted-variable bound of Cinelli and
Hazlett is the sensitivity row, and it is an oracle here \\
falsifier & the interval at one revision failing to cover at one wrong claim, or costing more than
discarding the knowledge. Neither happens in any cell \\
what it must not become & a continuity bound. Damage beyond the radius is heavy-tailed, so
``small knowledge error implies small bias'' is false and our own data falsifies it \\
carried by & T1 \\
\bottomrule
\end{tabular}
\end{center}

\subsection*{C5 \quad The tie: the validity half computes what a free screen computes}

A reviewer who knows the area will ask why an analyst needs a radius when they could leave one
claim out, close the graph again and look at whether the set changed. The answer is that at one
revision they could not tell the difference, because the two are the same policy, and at two
revisions they tie exactly when the screen is also run at two.

\begin{center}\footnotesize
\begin{tabular}{@{}p{2.6cm}p{13.75cm}@{}}
\toprule
rests on & the stop rule and the leave-one-out screen decide differently on 2 of 162 rows at one
wrong claim, both of the class the argument predicts. At two wrong claims the stop rule beats the
leave-one-out screen by 0.305 [0.182, 0.410] in coverage at the population and by
0.256 [0.149, 0.346] at a sample of one thousand, and ties the leave-two-out screen at
0.000 [0.000, 0.000]. The separation the design asked for exists only because that baseline stops
at depth one \\
the sentence this licenses & a screen answers yes or no at whatever depth the analyst happened to
choose. The radius says which depth was needed, and the coverage is won there. Quote one of the two
numbers above with its sample size attached; they are the same contrast measured twice \\
significance & this was pre-declared on 7 September, with the sentence for the null written before
the measurement, which is the reason it can be printed as a property rather than conceded as a
defeat. It is also the baseline the draft's ranking table does not have, and the single most
valuable thing our side can hand that table \\
nearest prior art & the screen is not in the literature; it is the thing a competent reader
invents in ten seconds, which is worse \\
the licensed sentence & the validity half of the report computes what a leave-$k$-out re-closure
screen computes, and the paper states this as a property of the measure. What the radius adds is
the depth, the named witness and the certificate that no failure lies nearer \\
carried by & T8, and one row of T1 \\
\bottomrule
\end{tabular}
\end{center}

\subsection*{C6 \quad The guarantee is expensive exactly where the knowledge matters}

On the problems where the knowledge carries identification, retracting a single true claim breaks
the committed set on nine problems in ten. The refusal rate is 0.910 on the informative rows with a
committed set, zero on the committed rows the class already identifies, and 0.639 over all
committed rows. The denominator decides the number, so the paper states the denominator.

\begin{center}\footnotesize
\begin{tabular}{@{}p{2.6cm}p{13.75cm}@{}}
\toprule
rests on & 134 informative committed rows over 20 networks in the semi-synthetic panel, 7 rows over
3 networks in the fitted panel; the prediction that the rate would stay below one half was
falsified \\
significance & this is the finding that stops the paper being a method paper. Knowledge buys
identifiability rather than a better set, and the certificate mostly says the conclusion is not
robust. That is a real diagnostic result and it reads as one only if we print it ourselves \\
falsifier & a population on which the refusal rate is low and the certificate still detects; the
informative stratum is not it \\
carried by & T8 \\
\bottomrule
\end{tabular}
\end{center}

\subsection*{C7 \quad Locality holds in the unit the analyst asserts, and not in hops}

Hop distance predicts the depth when the binding claim is incident to the treatment, and overstates
it otherwise. On the primary condition's 157 exact rows the hop reading equals the claim depth on
135, exceeds it by one on 19 and by two on 2, and falls below on 1. Reading above the claim depth
is what over-certification is.

\begin{center}\footnotesize
\begin{tabular}{@{}p{2.6cm}p{13.75cm}@{}}
\toprule
rests on & the agreement counts above, computed for this document; and the matched-budget negative
control, where retracting a claim far from the query is not inert on the semi-synthetic panel
(coverage $+0.044$ at one wrong claim) and is inert on the fitted panel \\
significance & locality pays in cost and in what is worth eliciting, and does not pay in width \\
what it must not become & ``knowledge at distance at least one cannot harm''. The 0/791 figure was
struck in August for having no provenance and is never repeated \\
carried by & T13, or the compact locality section \\
\bottomrule
\end{tabular}
\end{center}

\subsection*{C8 \quad Applicability, stated as a limit}

Of the nine author-declared queries in the corpus that can be certified, eight are unbroken by any
retraction down to depth three. Of the sampled population, the committed set breaks at the first
retraction on almost every row. The mechanism is the sampling rule: a pair is admitted only when
some retraction changes validity, so a robust query is rejected before it is measured.

\begin{center}\footnotesize
\begin{tabular}{@{}p{2.6cm}p{13.75cm}@{}}
\toprule
rests on & the observables-only frame, 80{,}132 candidate pairs over 33 networks, against 543 under
a gate that reads the true graph. Eight networks yield a frame where the old gate yielded no
instance at all \\
significance & the measured population is close to the complement of the queries the source papers
exist to answer. A reviewer who notices this before we state it has found our worst objection; a
reviewer who reads it in our own funnel has read a scope statement \\
carried by & T6, and it belongs in the funnel rather than in Table~1 \\
\bottomrule
\end{tabular}
\end{center}

\subsection*{C9 \quad Scope: the graph is the oracle's}

The licensed phrasing is that the radius is computable from the estimated CPDAG, the claims, the
covariance and the sample size. It is never that the conclusion is robust under an estimated CPDAG.
On the learned-graph panel the true graph extends the analyst's graph on none of the rows, and both
instruments over-certify on 14 of 38.

\subsection*{C10 \quad The supplier is an error process, and never the headline}

Measured commission and validity by condition, with a matched chance control: the model beats
chance in commission by 0.143 [0.006, 0.287] and is not demonstrated to beat it in validity,
0.082 [$-$0.020, 0.176]. Scale buys claims and accuracy and not validity: the share of committed
sets invalid at the truth is 0.298 at the largest supplier against 0.288 at the smallest. This is
methodology. It is ruled out of the Table~1 slot.

\subsection*{C11 \quad Ranking: the radius orders instances by survival better than the standard distances}

The radius ranks instances by whether the committed set survives corrupted knowledge better than
structural Hamming distance and than the number of claims, in most strata. Asserting more knowledge
is associated with less robustness. Structural Hamming distance is a poor predictor: knowledge can
be entirely correct and the set still break at the first subsequent error. One cell is a confirmed
null and two are unstable at 200 draws.

\begin{verdict}
\textbf{Not claimed, in any draft.} No continuity bound. No probability reading of the radius, and
no claim that it estimates how wrong an expert is. No ordering claim about suppliers. No use of the
word ``first''. No claim that knowledge at distance at least one cannot harm. No claim of
robustness under an estimated CPDAG. Retraction as a minimal correction set is cited rather than
named as ours.
\end{verdict}
